Read in Gauge Theory Gravity: the metric $g$ below is the derived metric of the gravitational field $(\hb, \Omega)$ on flat spacetime, the manifold is the invertible region, and the tractor connection is the flat-space $\Cl(2,4)$ gauge field of Theorem \claimref{synthesis.gtg-tractor-gauge-field-reproduces-tractor-connection}; the index computation below is its reading in the splitting of $g$ (Fact \claimref{synthesis.imported-gtg-dictionary}).

(1) is a computation with the bivector action $(a\wedge b)(v) = \ip{b}{v}a - \ip{a}{v}b$, the action of bivectors on vectors of \claimref{std-algebra.bivector-lie-algebra}, in the splitting $(\sigma, \mu, \rho)$ with $u_p$ the $\sigma$-direction and $X_p$ the $\rho$-direction. The bivector $u_p\wedge e_a$ maps the $\rho$-direction to $-e_a$ and $e_b$ to $g_{ab}\,u_p$, so $-(u_p\wedge e_a)$ contributes $(-\mu_a)$ to the $\sigma$-slot and $(+g_{ab}\rho)$ to the $\mu$-slot; and $(X_p\wedge \mathrm{P}(e_a))$ maps $u_p \mapsto -\mathrm{P}(e_a)$ and $e_b \mapsto \mathrm{P}_{ab}X_p$, contributing $(-\mathrm{P}_{ab}\sigma)$ to the $\mu$-slot and $(+\mathrm{P}_a{}^b\mu_b)$ to the $\rho$-slot. Summing reproduces the display of Fact \claimref{correspondence.imported-tractor-connection-local-twistor-transport}(i) entry by entry. The summand identifications follow: $X_p\wedge W_\C$ annihilates $X_p$, so it lies in the nilradical of the stabiliser of $X_p$ (Fact \claimref{correspondence.imported-stabilisers-flat-model} with $X_p$ in place of $\nInf$), while $u_p \wedge W_\C$ moves $X_p$: soldering.
(2) The grading eigenvalues are the two-line Fock computation recorded in Section \ref{sec:palatial} (Proposition \claimref{correspondence.clifford-side-ordering-shift}): $E = u_p\wedge X_p = -(X_p\wedge u_p)$ acts on the module (the Fock model of \claimref{correspondence.fock-model-complexified-ascended-algebra}) as $\tfrac12 c(u_p)c(X_p) - \tfrac12 = \tfrac12 - N$ with $N$ the occupation number of the $X_p$-mode (the two occupation numbers of a pair sum to one), so with eigenvalue $+\tfrac12$ on the states not containing the $X_p$-mode ($\Som$) and $-\tfrac12$ on those containing it ($\Spi$). For the summands: $[E, u_p\wedge w] = (u_p\wedge w)$ and $[E, X_p\wedge w] = -(X_p\wedge w)$ for $w \in W_\C$ ($E$-weights of $u_p, X_p, w$ are $+1, -1, 0$), so the soldering summand shifts the $E$-grading of spinors by $+1$ and the Schouten summand by $-1$; on a module whose $E$-spectrum is $\{\pm\tfrac12\}$, a weight-$(+1)$ operator maps the $(-\tfrac12)$-space to the $(+\tfrac12)$-space and annihilates the latter, and dually. The verification script confirms both summand actions on all four mixed directions in the Fock model.
(3) The block form is (1) and (2): $\nabla^g$ couples the Levi factor, which preserves both planes, and the two summands are the off-diagonal blocks. Equivariance of the summands: for $\ell$ in the Levi factor, $\ell$ fixes $u_p$ and $X_p$, so $[\rho(\ell), c(u_p)c(e)] = c(u_p)c(\ell e)$ and $[\rho(\ell), c(X_p)c(p)] = c(X_p)c(\ell p)$ by Proposition \claimref{correspondence.spin-representation}; thus $e_a \mapsto \sigma_a$ and $p_c := \mathrm{P}(e_c) \mapsto \tau_c$ are Levi-equivariant maps $W_\C \to \operatorname{Hom}(\Spi,\Som)$ and $W_\C\to\operatorname{Hom}(\Som,\Spi)$, i.e.\ elements of the one-dimensional spaces of Lemma \claimref{correspondence.chiral-planes-summand-spaces}(2), as are Penrose's two contractions; this gives the constants $c_1, c_2$. The composite: with $\rho(u_p\wedge e_a) = \tfrac12 c(u_p)c(e_a)$ and $\rho(X_p\wedge p_c) = \tfrac12 c(X_p)c(p_c)$ (both pairs orthogonal), for $\omega \in \Som$,
\[
\sigma_a\tau_c\,\omega = -\tfrac14\,c(u_p)c(e_a)c(X_p)c(p_c)\,\omega = \tfrac14\,c(u_p)c(X_p)\,c(e_a)c(p_c)\,\omega = \tfrac12\,c(e_a)c(p_c)\,\omega ,
\]
using $c(e_a)c(X_p) = -c(X_p)c(e_a)$, that $c(e_a)c(p_c)$ preserves $\Som$ (it does not involve the $X_p$-mode), and $c(u_p)c(X_p) = 2$ on $\Som$, which is Proposition \claimref{correspondence.clifford-side-ordering-shift} read on the $+\tfrac12$ eigenspace. The decomposition into $\tfrac12\ip{e_a}{p_c}$ and $\rho(e_a\wedge p_c)$ is the definition of $\rho$ on a non-orthogonal pair. The script checks the composite identity on all sixteen pairs of mixed directions. Substituting one summand of Fact \claimref{correspondence.imported-tractor-connection-local-twistor-transport}(ii) into the other gives Penrose's composite $I_7\,\varepsilon_A{}^B\,(I_7\,\mathrm{P}_{CC'DA'}\omega^D) = -\varepsilon_A{}^B\mathrm{P}_{CC'DA'}\omega^D$, his imaginary unit acting on the fibre $\CS$ as its scalar unit $I_7 = I_6e$, the central pseudoscalar of $\Cl(3,4) = \Cl(2,4)_\C$ (\claimref{std-algebra.complexification-is-the-odd-extension}), the identifications with $S^B$, $S_{B'}$ being linear over it. Curvature is not needed for the identification and is not used.
